\cotefile{82}
\scanpage{65}
à la section \deleted{\unsure{0|$\bullet$}} \inserted{nulle $s_0$} de \deleted{$\unsure{\mathbb P^1|\mathbb P}$} $X_0=\mathbb P^1$ (la section $\infty$, correspondant à la droite $\mathcal O_S e_2$, a comme stabilisateur
\[
\widetilde B'_0=\left\{\begin{pmatrix}\lambda&\unsure{0|b}\\ \mu&\lambda^{-1}\end{pmatrix}\ \middle|\ 
\begin{aligned}\lambda&\ \text{de }\unsure{\mathbb G_m|G_m}\\ \mu&\ \text{de }\mathbb G_a\end{aligned}\right\} \tag*{(11)}
\]

Comme l'action de $G_0$ sur $X_0$ \deleted{définit} section $s_0$ de $X_0$ définit
\[
\begin{aligned}g&\longmapsto g.s_0\\G_0&\longrightarrow X_0\end{aligned}\tag*{(12)}
\]
qui par passage au quotient définit un \uline{iso}
\[
G_0/B_0\longrightarrow X_0\ ,\tag*{(13)}
\]
\deleted{on trouve} compatible avec les op. de $G_0$, opérant par translation à g. sur \unsure{le} premier \deleted{\unsure{m}}\unsure{mb|b}, on trouve, en prenant l'image inverse sur $S$ et en tordant l'iso obtenu par le torseur $R$, l'iso.
\[
R\wedge X_{0S}\simeq R/B_{0S}\tag*{(14)}
\]
qui établit l'interprétation (7).

Le foncteur allant des \inserted{chacune des huit} catégories $C_{1S}\ldots$ vers $\operatorname{Tors}(\mathrm{GP}(1)_S)$ est le foncteur composé de
% archive sheet number: 64
\scanpage{66}
\[
\unsure{\mathcal X_0|\mathcal X|\mathfrak X_0}\longmapsto\underline{\operatorname{Isom}}(\mathcal X_0,\mathcal X)\tag*{(15)}
\]
vers \deleted{les} $\operatorname{Tors}(\underline{\operatorname{Aut}}(\mathcal X_{0S}))$, avec l'équivalence
\[
\operatorname{Tors}(\underline{\operatorname{Aut}}(\mathcal X_{0S}))\simeq\operatorname{Tors}(\mathrm{GP}(1)_S)\tag*{(16)}
\]
déduite de l'iso \deleted{\unsure{(7)}}
\[
\mathrm{GP}(1)_S\xrightarrow{\sim}\underline{\operatorname{Aut}}(\mathcal X_{0S})\bigl(\simeq\underline{\operatorname{Aut}}(\mathcal X_0)_S\bigr)\tag*{(17)}
\]
déduit de (7) par chgt de base.

\inserted{\unsure{Également} forme \unsure{particulière} du foncteur ``twist'' dans le cas de $C_{0S}\to C_{1S}$, cf ci-dessous}

\uline{IV} Le foncteur allant de la catégorie \deleted{$C_{1S}$ des formes de $\mathrm{GP}(1)_S$ vers des catégories} \deleted{catégories} $C_{0S}, C_{2S}$ etc vers la catégorie \inserted{$C_{1S}$} des formes de $\mathrm{GP}(1)$ s'explicite par
\[
\mathcal X\longrightarrow\underline{\operatorname{Aut}}(\mathcal X)\tag*{(18)}
\]

\deleted{Pour expliciter les foncteurs} Dans le cas \unsure{de|des} $C_{1S}'$ des formes $\widetilde G$ de $\mathrm{SL}(2)$, il y a aussi une description plus terre à terre
\[
\widetilde G\longrightarrow\widetilde G/\operatorname{Cent}(\widetilde G)\tag*{(19)}
\]
(\uline{NB} le centre de $\widetilde G$ est tjrs \deleted{\unsure{can}} isom à $\mu_2$ (de façon unique car $\operatorname{Aut}(\mu_2)=\overline1$

Pour le cas des $\underline{\operatorname{Quat}}(\mathcal O_S)$ $\to\operatorname{For}(G)\ldots)$
\[
M\longmapsto\check{\mathbb V}(M)^*/G_{mS}\tag*{(20)}
\]
% archive sheet number: 65
\scanpage{67}
Pour les foncteurs en sens inverse, allant de $C_{1S}=\operatorname{For}(\mathrm{GP}(1)_S)$ vers les autres catégories, on peut, soit passer par le torseur $R=\underline{\operatorname{Isom}}(G_{0S},G)$ puis tordre $\mathcal X_{0S}$
\[
G\unsure{\longmapsto|\vdash}\underline{\operatorname{Isom}}(G_{0S},G)\wedge_{G_{0S}}\mathcal X_{0S}\ ,\tag*{(21)}
\]
soit trouver, dans divers cas, des descriptions plus directes. Par exemple
\[
\begin{aligned}
G&\longmapsto\widetilde G=\text{Rev.\ simplement connexe}\\
 &\hspace{6em}\text{de $G$, défini comme}\\
 &\hspace{6em}\text{le schém. en groupes \unsure{ss|SS}}\\
 &\hspace{6em}\text{lisse à fibres \deleted{\unsure{géométr.|géom. connexes}} simplement}\\
 &\hspace{6em}\text{connexes, muni d'}\\
 &\hspace{6em}\text{une isog. $\widetilde G\to G$}\\
 &\hspace{5em}\text{-- lequel est défini de}\\
 &\hspace{6em}\text{façon \uline{unique} \ldots.}
\end{aligned}\tag*{(22)}
\]

Le foncteur vers les algèbres de Lie est
\[
\begin{aligned}
G&\longmapsto\mathfrak g=\underline{\operatorname{Lie}}(G)\\
\operatorname{For}(\mathrm{GP}(1)_S)&\xrightarrow{\unsure{\sim}}\operatorname{For}(\mathfrak{gp}(1)_S)
\end{aligned}\tag*{(23)}
\]
\[
\begin{aligned}
G&\longrightarrow\widetilde{\mathfrak g}=\underline{\operatorname{Lie}}(\widetilde G)\\
\operatorname{For}(\mathrm{GP}(1)_S)&\xrightarrow{\unsure{\sim}}\operatorname{For}(\mathfrak{sl}(2)_S)
\end{aligned}\tag*{(24)}
\]

Le foncteur
\[
G\longmapsto(E,q,\omega)
\]
peut se factoriser par $G\to\widetilde{\mathfrak g}$, cf plus bas.
% archive sheet number: 66
\scanpage{68}
Le foncteur
\[
G\longmapsto X
\]
s'explicite comme
\[
G\longmapsto\underline{\operatorname{Bor}}(G)\tag*{(2\unsure{5|4})}
\]
associant à $G$ le \deleted{\unsure{ens}} schém. de ses sous-groupes de Borel $B$ (i.e. \deleted{\unsure{tels}} des ss-gpes tels que \deleted{$(B,G)$} \inserted{$(G,B)$} \unsure{sont|soit} loc. isom. \deleted{\illegible{2}} (\inserted{\unsure{Zar}}, ou $\mathrm{fpqc}$) à \deleted{$\unsure{B|(G_0}$} $(G_{0S},B_{0S})$. On obtient en fait, si $G$ et $X$ se correspondent -- donc $G\simeq\underline{\operatorname{Aut}}(X)$ -- un isom. can.
\[
X\xrightarrow{\sim}\underline{\operatorname{Bor}}(G)\tag*{(2\unsure{6|5})}
\]
en associant à tte section \inserted{\unsure{locale|loc.}} de $X$ son stabilisateur dans $G$.

\uline{V} Les foncteurs vers les formes de $\mathrm{SL}(2)_S$ s'\deleted{\unsure{fac|f}}obtiennent soit en factorisant par $\operatorname{Tor}(\mathrm{GP}(1)_S)$
\[
\mathcal X\longmapsto R=\underline{\operatorname{Isom}}(\mathcal X_{0S},\mathcal X)\longmapsto R\times_{\mathrm{GP}(1)_S}\unsure{\widetilde G_0|\widetilde G_{0S}},\tag*{(27)}
\]
soit plus \unsure{naturellt|naturellement} en
\[
\mathcal X\longmapsto G=\underline{\operatorname{Aut}}_S(\mathcal X)\longrightarrow\widetilde G=\text{Rev. \unsure{simplemt|simpl. conn.} de $G$}\tag*{(28)}
\]
% archive sheet number: 67
\scanpage{69}
Dans le cas de \deleted{\unsure{Quat($S$)}} \deleted{Les foncteurs en sens inverse s'obtiennent}
\[
\operatorname{Quat}(\unsure{\mathcal O_S|S|\mathscr S})\longrightarrow\operatorname{For}(\mathrm{SL}(2)_S)
\]
on trouve
\[
M\longrightarrow\widetilde G\simeq\deleted{\unsure{\check{\mathbb V}|\check{\mathrm V}}(M)^*/\unsure{G_m}}\operatorname{Ker}\bigl(\check{\mathbb V}(M)^{\unsure{*|1}}\xrightarrow{\det}G_m\bigr)\tag*{(29)}
\]
\inserted{$=$ \uline{norme réduite}}

Pour les foncteurs de $\operatorname{For}(\mathrm{SL}(2)_S)$ vers les autres catégories, on peut soit passer par $\operatorname{Tor}(\mathrm{GP}(1))$ ou par $\operatorname{For}(\mathrm{GP}(1)_S)$, soit dans certains cas procéder directement, p.ex.
\[
\begin{aligned}
\operatorname{For}(\mathrm{SL}(2)_S)&\longrightarrow\operatorname{For}(\mathfrak{sl}(2)_S)\\
\widetilde G&\longrightarrow\widetilde{\mathfrak g}=\underline{\operatorname{Lie}}(\widetilde G)
\end{aligned}\tag*{(30)}
\]
\[
\widetilde G\longrightarrow X\simeq\underline{\operatorname{Bor}}(\widetilde G)\tag*{(31)}
\]

Pour la relation \deleted{\unsure{triple} entre $G,\widetilde G,\unsure{(E,q,\omega)|M,q},\mathfrak g,\widetilde{\mathfrak g}$} \inserted{\unsure{des ob}} \deleted{\unsure{on}} $G,\widetilde G,\mathfrak g,\widetilde{\mathfrak g},(E,q,\omega)$, c'est $M$ qui \unsure{semble|donne} peut être la meilleure clef pour les obtenir tous -- ainsi que $X$ et $(P,C)$, tant qu'à faire ! -- et pour comprendre les relations entre tous ces objets. On verra ainsi que $\mathfrak g$ et $\widetilde{\mathfrak g}$ sont des \unsure{modules duels|modules duaux} l'une de l'autre, que $E\simeq\widetilde{\mathfrak g}$ canoniquement i.e. qu'en \unsure{parti-|parti.}%
% archive sheet number: 68
\scanpage{70}
par $\widetilde{\mathfrak g}$, on verra \deleted{\unsure{des}} interprétations remarquables de $\widetilde{\mathfrak g}\to\mathfrak g$ et $\widetilde G\to G$ etc\ldots

\uline{VI} Soit donc $M$ une algèbre de quaternions sur $S$ -- i.e. une alg. loc. \inserted{(fpqc)} iso. à \deleted{$\unsure{\mathcal O}$} $\mathbb M_2(\mathcal O_S)$. On sait que
\[
\left\{\begin{aligned}
&\text{Centre de }M\deleted{\unsure{(\mathcal O_S)}}=\mathcal O_S.1\\
&(\text{Centre de }\deleted{\unsure{\check{\mathbb V}(M)^*|\mathrm V(M)^*}}M^*)=\mathcal O_S^*.1\quad\text{ou enc}\\
&(\text{Centre de }\mathcal L:=\check{\mathbb V}(M)^*)=G_m\subset\mathcal L(M)
\end{aligned}\right. \tag*{(32)}
\]

\deleted{\unsure{(33)}} qu'on a des hom ``\unsure{denom|norme} réduite'' et ``trace réduite''
\diagram{\begin{tikzcd}[ampersand replacement=\&,column sep=large]
\text{(33)}\&\check{\mathbb V}(M)\arrow[r,shift left=0.7ex,"\det"]\arrow[r,shift right=0.7ex,"\operatorname{tr}"']\&\deleted{\mathcal O_S}\mathbb E^1
\end{tikzcd}}
\margin{\uline{NB} $\det(uv)=\det u\det v$\quad $\left\{\begin{aligned}&\operatorname{tr}\text{ est $\mathcal O_S$-linéaire}\\&\det\text{ est quadratique}\end{aligned}\right.$\quad $\det(u+\lambda.1)=\deleted{\det u\det\unsure{\lambda}}\lambda^2+\lambda\operatorname{Tr}u+\det u$\quad L'appl. tangente à $\det$ en la section $1$ est $\operatorname{tr}$}
définis par des conditions évidentes (par \deleted{\unsure{descente}} \inserted{localisation} : partir de $\det$ et $\operatorname{tr}$ sur $\mathbb M_2(\mathcal O_S)$, \deleted{\unsure{puis}} qui sont inv. par autom)

On obtient des suites exactes
\[
1\longrightarrow G_m\xrightarrow{\substack{\text{incl du}\\\text{centre}}}\mathcal L(M)\longrightarrow G\longrightarrow 1\tag*{(34)}
\]
\[
1\longrightarrow\widetilde G\longrightarrow\mathcal L(M)\xrightarrow{\det}G_m\longrightarrow1\tag*{(35)}
\]
et sur les algèbres de Lie
\[
0\longrightarrow\mathcal O_S\xrightarrow{\lambda\mapsto\lambda.1}M\longrightarrow\mathfrak g\longrightarrow0\tag*{(34$'$)}
\]
\[
0\longrightarrow\widetilde{\mathfrak g}\longrightarrow M\xrightarrow{\operatorname{tr}}\mathcal O_S\longrightarrow0\tag*{(35$'$)}
\]
% archive sheet number: 69
\scanpage{71}
Ces suites exactes déterminent $G,\widetilde G,\mathfrak g,\widetilde{\mathfrak g}$ chacun directement en termes de $\mathcal L(M)$ et $\det$ (pour les groupes) et $M$ et $\operatorname{tr}$ (pour les algèbres de Lie). Bien entendu, $M$ est considéré comme algèbre de Lie \unsure{grâce} à la forme \unsure{trace}.

Les hom canoniques
\[
\left\{\begin{aligned}
\widetilde G&\longrightarrow G\\
\widetilde{\mathfrak g}&\xrightarrow{f}\mathfrak g
\end{aligned}\right.\tag*{(36}
\]
sont les composés
\[
\widetilde G\overset{\mathrm{incl.}}{\lhook\joinrel\longrightarrow}\mathcal L(M)\xrightarrow{\mathrm{proj}}G\tag*{(37)}
\]
\[
\widetilde{\mathfrak g}\overset{\mathrm{incl}}{\lhook\joinrel\longrightarrow}M\xrightarrow{\mathrm{proj}}\mathfrak g\tag*{(38)}
\]

Considérons sur $M$ la forme \deleted{quadratique} \inserted{bilinéaire} \deleted{\unsure{obtenue}}.
\[
\deleted{\unsure{\Phi|q}}\Phi(uv)=\operatorname{Tr}(uv)\tag*{(39)}
\]
On voit alors que $\varphi$ est symétrique et \unsure{$f$|$\varphi$} non dégénérée en chaque fibre, permettant donc d'identifier $M$ à son propre dual
\[
M\simeq M^\vee\tag*{(40)}
\]
L'orthogonale de $\mathcal O_S$ est $\widetilde{\mathfrak g}$, et on trouve ainsi un accouplement dualisant
% archive sheet number: 70
\scanpage{72}
\[
\widetilde{\mathfrak g}\times\mathfrak g\longrightarrow\mathcal O_S\tag*{(41)}
\]
donnant donc des isom de Modules :
\[
\widetilde{\mathfrak g}\simeq\overset{\deleted{\unsure{\sim}}}{\mathfrak g}^\vee,\quad\mathfrak g\simeq\widetilde{\mathfrak g}^{\vee}\tag*{(42)}
\]

La \deleted{\unsure{a sur}} forme quadratique \inserted{$Q(u)=$} $-\det u$ sur $M$ est lisse, et la forme bilinéaire associée est
\[
\left\{\begin{aligned}
\Phi_1(\unsure{uv|u,v})&=-\operatorname{Tr}u\operatorname{Tr}v+\operatorname{Tr}(uv)=-\operatorname{Tr}u\operatorname{Tr}v+\Phi(u,v)\\
\text{donc}\quad-\det(u\unsure{+}v)+\det u+\det v&=\Phi_1(u,v)\\
&=-\operatorname{Tr}u\operatorname{Tr}v+\operatorname{Tr}(uv)
\end{aligned}\right.\tag*{(43)}
\]
qui est donc également dualisante. Sur $\widetilde{\mathfrak g}$ elle coïncide avec $\varphi$, et \deleted{\unsure{elle a}} comme $\operatorname{Tr}(1)=2$, on a
\[
\Phi_1(1,u)=-\Phi(1,u)=-\operatorname{Tr}u
\]
donc l'ortho. de $\mathcal O_S.1$ pour $\varphi_1$ est encore $\widetilde{\mathfrak g}$. L'accouplement \deleted{\unsure{le}} \inserted{$\mathfrak g\times\widetilde{\mathfrak g}\to\mathcal O_S$} déduit de $\varphi_1$ par passage au quotient \deleted{\unsure{diffère par le signe de}} \inserted{est le $=$ que celui} (41) défini par $\Phi$ ; \deleted{\unsure{puis que}} \inserted{D'ailleurs} \deleted{\unsure{la la r}} $\Phi$ et $\Phi_1$ ont $=$ restriction à $\widetilde{\mathfrak g}$, soit $\varphi$
\[
\left\{\begin{aligned}
\varphi&=\Phi\mid\widetilde{\mathfrak g}\times\widetilde{\mathfrak g}
=\Phi_1\mid\widetilde{\mathfrak g}\times\widetilde{\mathfrak g}
:\widetilde{\mathfrak g}\times\widetilde{\mathfrak g}\to\mathcal O_S\\
\varphi(uv)&=\operatorname{Tr}uv
\end{aligned}\right.\tag*{(42)}
\]
% archive sheet number: 71
\scanpage{73}
Considérons la forme quadratique $\unsure{q_0|q}$ sur $\widetilde{\mathfrak g}$
\[
\left\{\begin{aligned}
q&=Q\mid\widetilde{\mathfrak g}:\widetilde{\mathfrak g}\longrightarrow\mathcal O_S\\
q(u)&=-\det u
\end{aligned}\right.\tag*{(43)}
\]
alors \deleted{$\unsure{q}$} $\varphi$ est donc associée à la forme $q$
\[
\left\{\begin{aligned}
q(x+y)&=q(x)\unsure{{}\,|+}q(y)+\varphi(x,y),\\
\varphi(u,u)&=2q(u)
\end{aligned}\right.\tag*{(44)}
\]
d'ailleurs $q$ est lisse -- $\varphi$ est dualisante sauf en car. $2$.

[La relation avec Killing est
\[
\left\{\begin{aligned}
\operatorname{Tr}(\operatorname{ad}_{\widetilde{\mathfrak g}}(u)\operatorname{ad}_{\widetilde{\mathfrak g}}(v))&=4\varphi(u,v)\\
\operatorname{Tr}(\operatorname{ad}_{\widetilde{\mathfrak g}}(u)^2)&=8q(u)
\end{aligned}\right.\tag*{(45)}
\]
pour $u,v$ sections de $\widetilde{\mathfrak g}$. \enspace]

On définit aussi une base
\[
\omega\in\Gamma\underline{\operatorname{\text{dét}}}\,\widetilde{\mathfrak g}\simeq\Gamma\underline{\operatorname{\text{dét}}}\,M\tag*{(46)}
\]
(le dernier iso provient de la suite exacte (35$'$) \unsure{en passant au dét}) compte tenu de l'isom local
\[
M\simeq\underline{\operatorname{End}}(V)\simeq\check V\otimes V
\]
et de l'isomorphisme canonique
\[
\begin{aligned}
\det(\check V\otimes V)&\simeq\det\check V\otimes\det V\simeq(\underline{\det}V)^{-1}\otimes\underline{\det}V\\
&\simeq\mathcal O_S
\end{aligned}\tag*{(47)}
\]
% archive sheet number: 72
\scanpage{74}
qui est évidemment stable par $\mathrm{GL}(V)$ donc par $\mathrm{GP}(V)$ -- donc se descend en une base $\omega$ de $\underline{\det}M$, d'où de $\det\widetilde{\mathfrak g}$. Il est vrai que dans l'iso (46) il y a une convention de signe à faire -- à défaut, il y a l'ambiguïté des signes (qu'on pouvait définir une \deleted{\unsure{telle}} $\omega$ canoniquement modulo cette seule ambiguïté était évident, car sur $\mathbb Z$ tout Module inversible a \unsure{exactt|exactement} deux bases, différant par le signe -- et l'invariance par \deleted{$\unsure{\mathrm{GP}1|\mathrm{GP}(1)}$} $G_0$ est claire. Mais $\omega^{\otimes2}$ ne dépend plus de cette convention, ce qui donne un sens dénué de tte ambiguïté à
\[
\underset{\substack{\text{discriminant}\\\text{divisé}}}{\underline{\delta'(q)}}=-\omega^{\otimes2}\tag*{(48)}
\]
d'où
\[
\delta(\varphi)=\delta(q)=-2\omega^{\otimes2}\tag*{(49)}
\]
% archive sheet number: 73
\scanpage{75}
On \deleted{posera}
\[
(E,q,\omega)=(\widetilde{\mathfrak g},q,\omega)\qquad\text{i.e}\quad E=\widetilde{\mathfrak g}\tag*{(50)}
\]

\deleted{On peut aussi définir,}

\uline{NB} L'hom \deleted{\unsure{$f$ de l|fondamental}}
\[
f:\widetilde{\mathfrak g}\longrightarrow\mathfrak g\tag*{(51)}
\]
de (36) ``n'est autre'' que l'hom $\widetilde{\mathfrak g}\to\overset{\deleted{\unsure{\sim}}}{\widetilde{\mathfrak g}}^\vee$ associé à la forme bilinéaire $\varphi$, compte tenu de l'iso $\widetilde{\mathfrak g}^\vee\simeq\mathfrak g$. Ceci signifie en effet que
\[
\begin{aligned}
\langle\overset{\deleted{\unsure{w|-}}}{u},f(v)\rangle&=\varphi(u,v)\qquad u,v\in\Gamma(S,\widetilde{\mathfrak g})\\
\shortparallel\qquad&\quad\shortparallel\\
\operatorname{Tr}uv&\quad\operatorname{Tr}uv
\end{aligned}
\]

Ceci dit, les lois de crochet sur $\widetilde{\mathfrak g}$ et sur son dual $\mathfrak g$ sont celles déduites canoniquement par ailleurs des couples $(\varphi,\omega)$ d'une forme bil. sym de $\widetilde{\mathfrak g}$ et d'une base $\omega$ de $\det\widetilde{\mathfrak g}$ (c'est ici que la convention de signes pour $\omega$ importe -- on prendra
\[
\left|\overline{\omega=\widetilde X\wedge\widetilde H\wedge\widetilde Y
=-\underset{\substack{\shortparallel\\e_1}}{\widetilde X}\wedge\underset{\substack{\shortparallel\\e_2}}{\widetilde Y}\wedge\underset{\substack{\shortparallel\\e_3}}{\widetilde H}\ .}\right|\tag*{(52)}
\]
\margin{Sans garantie -- il \deleted{faut} vérifier \unsure{le texte|les textes} \ldots}
% slanted
\deleted{On}
% archive sheet number: 74
\scanpage{76}
\inserted{On} peut aussi expliciter dans notre contexte actuel l'hom
\[
f':\mathfrak g\longrightarrow\widetilde{\mathfrak g}\tag*{(53)}
\]
(s'identifiant à $\Lambda^2f$) tel que
\[
\left\{\begin{aligned}
ff'&=-2\operatorname{id}_{\mathfrak g}\\
f'f&=-2\operatorname{id}_{\widetilde{\mathfrak g}}
\end{aligned}\right.\tag*{(54)}
\]
en posant
\[
f'(\bar u)=(\operatorname{Tr}u)1-2u\tag*{(55)}
\]
où $u$ est une section loc de $M$, et $\bar u$ son image par $M\to\mathfrak g$ ; comme $\operatorname{Tr}1=2$ cette application \deleted{\unsure{en $u$ est}} $M\to M$ est bien à val. dans $\widetilde{\mathfrak g}$, et nulle sur $\mathcal O_S.1$ donc se factorise par $\mathfrak g$, enfin si $u$ est dans $\unsure{\widetilde{\mathfrak g}|\mathfrak g}$ on trouve $f'f(u)=-2u$, enfin $ff'(\bar u)=\deleted{\unsure{-2f(u)}}\overline{f'(\bar u)}=-2\bar u$ OK.

Il reste à voir comment s'explicitent en termes de $M$ les objets $X$, et $P\supset\deleted{\unsure{C}}X$. Bien sûr, on aura
\[
\left\{\begin{aligned}
P&=\check{\mathbb P}(\widetilde{\mathfrak g})\simeq\mathbb P(\mathfrak g)\\
X&=\mathbb V(q)\subset P
\end{aligned}\right.\tag*{(56)}
\]
\margin{\deleted{\uline{NB} les sections de $X$ sont \unsure{en corresp biuniv|aussi les suiv.} :}}
% slanted
% archive sheet number: 75
\scanpage{77}
La relation \deleted{\unsure{avec}} entre cette description du fibré en droites projectives $X$, et la description comme $\underline{\operatorname{Bor}}(G)$ ou $\underline{\operatorname{Bor}}(\widetilde G)$, est explicitée dans la

\uline{Proposition} Les ss-groupes de Borel \inserted{$B$} de $G$ correspondent \unsure{biunivqt|biunivoquemt|biunivoquement} par image inverse à \deleted{ceux} \inserted{$B_1$} de $\mathcal L=\mathbb V(M)^\times$ ou \inserted{$\widetilde B$} de $\widetilde G$. Un ss-groupe de Borel est connu quand on connaît son algèbre de Lie $L$ (qui est un ss fibré vectoriel de corang $1$ dans $\mathfrak g$), ou celle de $L_1$ de $\mathcal L$ qui lui correspond (qui est un ss-fibré \deleted{de rang \unsure{2|3}} \inserted{vectoriel} de corang $1$) image inverse de $L_1$ dans $M=\underline{\operatorname{Lie}}(\mathcal L)$, \deleted{ou celle} ou \uuline{en car.} $\ne2$ quand on connaît celle $\widetilde L$ de $\widetilde B$ dans $\widetilde{\mathfrak g}$, qui n'est autre que $\widetilde L\cap\widetilde{\mathfrak g}$ -- (et pour cause c'est \unsure{pareil|permis}, puisqu'alors $f:\widetilde{\mathfrak g}\xrightarrow{\sim}\mathfrak g$ !). \deleted{Pour} \inserted{En termes} \deleted{\unsure{pour un ss-fibré vectoriel}} de $\widetilde L$ on aura
\[
\widetilde B=\check{\mathbb V}(\widetilde L)\cap\underbrace{\mathbb V(M)}_{\mathcal L}^{*}\qquad(B=\widetilde B/G_m)\tag*{(57)}
\]
% archive sheet number: 76
\scanpage{78}
Pour qu'un sous-fibré vectoriel \inserted{$L$} de $\mathfrak g$ de corang $1$ soit de la forme $\underline{\operatorname{Lie}}(B)$, il f. et s. que le fibré \deleted{$\unsure{L}$} \inserted{$L^\perp$} de rang $1$ orthogonal \deleted{$\unsure{L^\perp}$} dans $\widetilde{\mathfrak g}$ satisfasse
\margin{\uline{Question} suffit-il que $L$ soit une ss-Algèbre de Lie i.e. stable par crochet ??}
% slanted
\[
(q\mid L^\perp)=0
\]
\margin{\uline{NB} Il \unsure{revient|reprends} au $=$ de dire que $L^\perp$ \deleted{\unsure{soit}} \inserted{est} formé d'él. de \uline{carré nul} de $M$}
% slanted
(et ceci, bien sûr, caractérise $q$ \inserted{sur $\widetilde{\mathfrak g}$ (défini comme $\mathfrak g^\vee$)} à scalaire inversible près \inserted{en termes de $G$} -- donc au signe près si on veut qu'il soit défini sur tte base de façon compatible avec chgt de base)

\deleted{\unsure{Hyp|NB}} \uline{NB} Plaçons-nous \uline{en car $2$} Alors
\[
\mathcal O_S1\subset\widetilde{\mathfrak g}\tag*{(58)}
\]
et le \deleted{\unsure{l'application|application}} morphisme
\[
X\longrightarrow\deleted{\unsure{\mathrm{Gr}}}\unsure{P|\mathbb P}(\widetilde{\mathfrak g})\tag*{(59)}
\]
allant des Borels de $\widetilde{\mathfrak g}$ vers leurs algèbres de Lie, \unsure{considéré|considérés|considérées} comme sous fibré vect. de corang $1$, est celui qui à tte section de $X$, identifiée à une \inserted{fibré en} \unsure{droites|droite} \inserted{$L$} dans $\widetilde{\mathfrak g}$ satisfaisant $q\mid L=0$, associe son orthogonal dans $\widetilde{\mathfrak g}$ pour $\varphi$. \deleted{L'image est donc les} \inserted{Comme} l'orthogonal contient $\mathcal O_S.1$, (qui est le noyau de $\varphi$ \unsure{\}|,} i.e. de $f$), \deleted{est} le morphisme précédent se factorise par
% archive sheet number: 77
\scanpage{79}
$\mathbb P(\widetilde{\mathfrak g}/\mathcal O_S)$, \deleted{\unsure{et}} d'où
\[
X\longrightarrow\mathbb P(\widetilde{\mathfrak g}/\mathcal O_S)\quad
\begin{aligned}&\text{qui est lui-$=$ un}\\&\text{fibré en droites projectives}\end{aligned}\tag*{(\unsure{69|60|59})}
\]
Si \inserted{(59) donc} ce morphisme \inserted{(60)} était encore un monomorphisme, il \deleted{aurait} résulterait que (60) est $=$ un isomorphisme, donc $X$ serait ``banal'' du pt de vue Brauer -- or il n'est pas vrai qu'en car $2$ les fibrés en droites projectives soient tjrs banals !

En fait, on trouve ici un iso. canonique
\[
\mathbb P(\widetilde{\mathfrak g}/\mathcal O_S)\simeq X^{(2)}\;\raisebox{.65ex}{\(\scriptstyle\longrightarrow\text{(frobénisé)}\)}.\tag*{(61)}
\]
et (60) n'est autre que le morphisme canonique
\[
X\longrightarrow X^{(2)}
\]
qui est un épimorphisme radiciel (et \deleted{plat} \inserted{fini loc. libre} de rang $2^d=2$), mais en effet nullement un monomorphisme \ldots.
% archive sheet number: 78
